% L6c lecture companion. Notebook hashes and coverage are recorded in README.md.
% Style files and Cornell seal are copied unchanged from the approved L6a deck.
\documentclass[aspectratio=169,11pt]{beamer}
\usepackage{vnslides}
\usepackage{tabularx}
\renewcommand{\vnlecturecode}{L6c}
\newcolumntype{Y}{>{\raggedright\arraybackslash}X}
\newcommand{\lecturelink}{https://github.com/varnerlab/CHEME-5800-CourseRepository-Fall-2026/blob/main/weeks/week-06/L6c/CHEME-5800-L6c-Lecture-GeneralIterativeMethod-Fall-2026.ipynb}
\newcommand{\examplelink}{https://github.com/varnerlab/CHEME-5800-CourseRepository-Fall-2026/blob/main/weeks/week-06/L6c/CHEME-5800-L6c-Example-FunWithIterativeSolvers-Fall-2026.ipynb}
\newcommand{\jacobilink}{https://github.com/varnerlab/CHEME-5800-CourseRepository-Fall-2026/blob/main/weeks/week-06/L6c/CHEME-5800-L6c-Algorithm-JacobiMethod-Fall-2026.ipynb}
\newcommand{\gslink}{https://github.com/varnerlab/CHEME-5800-CourseRepository-Fall-2026/blob/main/weeks/week-06/L6c/CHEME-5800-L6c-Algorithm-GaussSeidel-Fall-2026.ipynb}
\newcommand{\sorlink}{https://github.com/varnerlab/CHEME-5800-CourseRepository-Fall-2026/blob/main/weeks/week-06/L6c/CHEME-5800-L6c-Algorithm-SOR-Fall-2026.ipynb}
\newcommand{\convergencelink}{https://github.com/varnerlab/CHEME-5800-CourseRepository-Fall-2026/blob/main/weeks/week-06/L6c/CHEME-5800-L6c-Advanced-Convergence-IterativeMethods-Fall-2026.ipynb}
\newcommand{\lablink}{https://github.com/varnerlab/CHEME-5800-CourseRepository-Fall-2026/blob/main/weeks/week-06/L6d/CHEME-5800-L6d-Lab-IterativeLinearSolvers-Fall-2026.ipynb}
\newcommand{\lecturesource}[1]{\note{[Sources]\par
  \href{\lecturelink}{L6c lecture: Iterative Methods for Solving Linear Algebraic Equations.}\par
  #1\par[/Sources]}}

\title{\texorpdfstring{{\vnheavy L6c}}{L6c}: Iterative Methods for Solving Linear Algebraic Equations}
\subtitle{\texorpdfstring{Residual corrections, convergence, and matrix splittings\\[3pt]CHEME 5800 Cornell University}{Residual corrections, convergence, and matrix splittings; CHEME 5800 Cornell University}}
\author{Copyright Jeffrey Varner 2026. All Rights Reserved}

\begin{document}
\special{pdf:minorversion 7}
\vntitlepage

\begin{frame}{Objectives for Today}
  How can repeated updates solve a linear system, and when do the
  approximations approach the solution?

  \vspace{0.4em}
  {\large\textbf{\ink{Today's concepts}}}
  \vspace{0.2em}
  \begin{itemize}\setlength{\itemsep}{0.7em}
    \item \hilite{Construct an iterative method}: derive a correction from
      a matrix splitting and distinguish satisfying a tolerance from reaching an iteration limit.
    \item \hilite{Assess convergence}: use the spectral radius, diagonal
      dominance, and an error bound to determine whether the iterates approach the solution.
    \item \hilite{Compare specific methods}: explain the updates and
      computational work of Jacobi, Gauss–Seidel, and relaxation.
  \end{itemize}
  \slidenote{Lecture:}{\href{\lecturelink}{Iterative Methods for Solving Linear Algebraic Equations}}
  \lecturesource{Introduction and three Learning Objectives.}
\end{frame}

\begin{frame}{Lecture and Example}
  \textbf{\ink{Lecture:}}
  \href{\lecturelink}{Iterative Methods for Solving Linear Algebraic Equations}

  Derive residual corrections, establish convergence conditions, and compare
  diagonal and triangular matrix splittings.

  \vspace{1.1em}
  \textbf{\ink{Example:}}
  \href{\examplelink}{Fun with Iterative Methods}

  Solve the same system with an iterative method and a direct solver.
  Check the computed solutions, then compare the solvers' runtimes and memory allocations.

  \vspace{0.8em}
  Repeated corrections offer an alternative when a direct solution requires
  too much computation or memory.
  \lecturesource{Examples; motivation in the Introduction.}
\end{frame}

\begin{frame}{General Iterative Method}
  Let $\mathbf A\in\mathbb R^{n\times n}$ be nonsingular, let
  $\mathbf x\in\mathbb R^n$ contain the unknowns, and let
  $\mathbf b\in\mathbb R^n$ be the right-hand side. We seek the unique solution of:
  \[
    \mathbf A\mathbf x=\mathbf b.
  \]
  An \hilite{iterate} $\mathbf x^{(k)}$ is the approximation after $k$ updates,
  starting from the initial guess $\mathbf x^{(0)}$. The \hilite{residual} at this iterate is:
  \[
    \mathbf r^{(k)}=\mathbf b-\mathbf A\mathbf x^{(k)}.
  \]
  Each component of the residual measures the mismatch in one equation.
  If the residual is zero, the iterate is the exact solution.

  \vspace{0.6em}
  The \hilite{solution error} is the difference between the approximation and
  the exact solution. A small residual alone does not guarantee a small solution error.
  \lecturesource{General iterative method: system, residual, and accuracy interpretation.}
\end{frame}

\begin{frame}{A Matrix Splitting Defines the Correction}
  Choose a nonsingular matrix $\mathbf M$ whose systems are inexpensive to
  solve. A \hilite{matrix splitting} writes $\mathbf A$ as:
  \[
    \mathbf A=\mathbf M-\mathbf N,
    \qquad \mathbf N=\mathbf M-\mathbf A.
  \]
  A \hilite{correction} $\mathbf d^{(k)}$ is the vector added to the current
  approximation. We compute the correction and update the approximation using:
  \[
    \underbrace{\mathbf M\mathbf d^{(k)}=\mathbf r^{(k)}}_{\text{solve for the correction}},
    \qquad
    \underbrace{\mathbf x^{(k+1)}=\mathbf x^{(k)}+\mathbf d^{(k)}}_{\text{apply the correction}}.
  \]
  A correction need not reduce the residual norm or the solution-error norm.

  \vspace{0.6em}
  A fixed splitting defines a \hilite{stationary iterative method}.
  Choosing a diagonal or triangular $\mathbf M$ simplifies the correction solve.
  \lecturesource{General iterative method: splitting and correction solve.}
\end{frame}

\begin{frame}{General Iterative Method: Algorithm}
  \textbf{Initialize:} Choose $\mathbf x^{(0)}$, an absolute tolerance $\epsilon>0$
  for the Euclidean residual norm, and a nonnegative integer update limit $\texttt{maxiter}$.
  Set $k\gets0$ and $\texttt{converged}\gets\texttt{false}$.

  \vspace{0.4em}
  While not $\texttt{converged}$ \textbf{do}:
  \begin{enumerate}\setlength{\itemsep}{0.15em}
    \item Calculate $\mathbf r^{(k)}\gets\mathbf b-\mathbf A\mathbf x^{(k)}$.
    \item Check whether to stop:
      \begin{itemize}\setlength{\itemsep}{0.1em}
        \item If $\|\mathbf r^{(k)}\|_2<\epsilon$, \textbf{then}: return
          $\mathbf x^{(k)}$ with $\texttt{converged}=\texttt{true}$.
        \item Otherwise, if $k\ge\texttt{maxiter}$, \textbf{then}: return
          $\mathbf x^{(k)}$ with $\texttt{converged}=\texttt{false}$ and a warning.
      \end{itemize}
    \item Solve $\mathbf M\mathbf d^{(k)}=\mathbf r^{(k)}$.
    \item Update $\mathbf x^{(k+1)}\gets\mathbf x^{(k)}+\mathbf d^{(k)}$.
    \item Increment $k\gets k+1$.
  \end{enumerate}
  \vspace{0.3em}
  The counter $k$ records completed updates. The \texttt{converged} flag reports
  whether the residual tolerance was met. Both return statements terminate the algorithm.
  \lecturesource{General iterative method: pseudocode. Preserve the 2025 Initialize/While/numbered-step/nested-conditional format.}
\end{frame}

\begin{frame}{Update Direction: The Stationary Form}
  Substitute the splitting into the system, then use the current approximation
  on the right-hand side:
  \[
    \begin{aligned}
      (\mathbf M-\mathbf N)\mathbf x&=\mathbf b,\\
      \mathbf M\mathbf x&=\mathbf b+\mathbf N\mathbf x,\\[3pt]
      \mathbf M\mathbf x^{(k+1)}&=\mathbf b+\mathbf N\mathbf x^{(k)}.
    \end{aligned}
  \]
  Define the iteration matrix $\mathbf G$ and constant vector $\mathbf c$:
  \[
    \mathbf x^{(k+1)}=
    \underbrace{\mathbf M^{-1}\mathbf N}_{\mathbf G}\mathbf x^{(k)}
    +\underbrace{\mathbf M^{-1}\mathbf b}_{\mathbf c}.
  \]
  The matrix $\mathbf G$ and vector $\mathbf c$ remain fixed as $k$ increases.
  We can also express this update as a correction added to the current approximation.
  \lecturesource{Update direction: splitting identity and stationary iteration.}
\end{frame}

\begin{frame}{Update Direction: A Residual Correction}
  With $\mathbf I$ the identity matrix, substitute
  $\mathbf N=\mathbf M-\mathbf A$ into the stationary update:
  \[
    \begin{aligned}
      \mathbf x^{(k+1)}
      &=\mathbf M^{-1}\mathbf b+(\mathbf I-\mathbf M^{-1}\mathbf A)\mathbf x^{(k)}\\[3pt]
      &=\mathbf x^{(k)}+\mathbf M^{-1}
        \underbrace{(\mathbf b-\mathbf A\mathbf x^{(k)})}_{\text{residual }\mathbf r^{(k)}}\\[3pt]
      &=\mathbf x^{(k)}+\underbrace{\mathbf M^{-1}\mathbf r^{(k)}}_{\text{correction }\mathbf d^{(k)}}.
    \end{aligned}
  \]
  Thus, the correction solves $\mathbf M\mathbf d^{(k)}=\mathbf r^{(k)}$.

  \vspace{0.6em}
  The exact solution is a \hilite{fixed point}: its residual is zero, so
  the update leaves the solution unchanged. We must still determine whether
  iterates starting from other initial guesses approach this fixed point.
  \lecturesource{Update direction: residual derivation and stationary-iteration callout.}
\end{frame}

\begin{frame}{Convergence: Propagation of the Error}
  Let $\mathbf x^\star$ be the exact solution and define the solution error
  by $\mathbf e^{(k)}=\mathbf x^{(k)}-\mathbf x^\star$.
  Subtract the fixed-point equation from the update:
  \[
    \begin{aligned}
      \mathbf e^{(k+1)}
      &=(\mathbf G\mathbf x^{(k)}+\mathbf c)
       -(\mathbf G\mathbf x^\star+\mathbf c)\\
      &=\mathbf G\mathbf e^{(k)},\\[5pt]
      \mathbf e^{(k)}&=\mathbf G^k\mathbf e^{(0)}.
    \end{aligned}
  \]
  The residual and the solution error are related by $\mathbf r^{(k)}=-\mathbf A\mathbf e^{(k)}$.

  \vspace{0.7em}
  \hilite{Convergence} means $\mathbf e^{(k)}\to\mathbf0$ as $k\to\infty$.
  The iteration converges from every initial guess if and only if $\mathbf G^k\to\mathbf0$.
  \lecturesource{Convergence: error recurrence. \href{\convergencelink}{Advanced convergence companion.}}
\end{frame}

\begin{frame}{The Spectral-Radius Convergence Condition}
  Let $\lambda_i$ denote the eigenvalues of $\mathbf G$.
  The \hilite{spectral radius} of $\mathbf G$ is the largest eigenvalue magnitude:
  \[
    \rho(\mathbf G)=\max_i|\lambda_i|.
  \]
  The stationary iteration converges to the solution from every initial
  guess if and only if:
  \[
    \boxed{\rho(\mathbf G)<1.}
  \]
  The condition $\rho(\mathbf G)<1$ guarantees that the error approaches zero.
  The error norm may nevertheless increase during some updates.
  \slidenote{Derivation:}{\href{\convergencelink}{Error propagation and the spectral-radius criterion}}
  \lecturesource{Convergence: spectral-radius condition. Advanced companion: necessity, sufficiency, and Jordan blocks.}
\end{frame}

\begin{frame}{Convergence Need Not Be Monotone}
  Take the following iteration matrix and initial error:
  \[
    \mathbf G=\begin{bmatrix}1/2&2\\0&1/2\end{bmatrix},
    \qquad\mathbf e^{(0)}=\begin{bmatrix}0\\1\end{bmatrix}.
  \]
  Both eigenvalues of $\mathbf G$ are $1/2$, but the first update increases the error norm:
  \[
    \mathbf e^{(1)}=\begin{bmatrix}2\\1/2\end{bmatrix},
    \qquad \|\mathbf e^{(1)}\|_2=\frac{\sqrt{17}}2\approx2.06>1.
  \]
  Repeated multiplication nevertheless gives:
  \[
    \mathbf G^k=2^{-k}\begin{bmatrix}1&4k\\0&1\end{bmatrix}\longrightarrow\mathbf0.
  \]
  The factor $2^{-k}$ decays faster than $k$ grows, so the error approaches zero.
  \lecturesource{Convergence: eventual versus monotone decay. Exact example from the advanced convergence companion, Convergence need not be monotone.}
\end{frame}

\begin{frame}{Diagonal Dominance}
  Write $a_{ij}$ for entry $(i,j)$ of $\mathbf A$.
  \hilite{Strict diagonal dominance by rows} requires:
  \[
    |a_{ii}|>\sum_{j\ne i}|a_{ij}|,
    \qquad i=1,\ldots,n.
  \]
  The magnitude of each diagonal entry exceeds the sum of the magnitudes
  of the other entries in that row.

  \vspace{0.8em}
  \begin{itemize}\setlength{\itemsep}{0.7em}
    \item This condition guarantees convergence of \hilite{Jacobi}
      and \hilite{Gauss–Seidel} from every initial guess.
    \item The inequality must hold in \hilite{every row}.
    \item Strict diagonal dominance is sufficient for convergence, but it is not necessary.
  \end{itemize}
  \lecturesource{Convergence: Diagonal dominance.}
\end{frame}

\begin{frame}{Why Diagonal Dominance Works for Jacobi}
  Assume that $\mathbf A$ is strictly diagonally dominant by rows.
  Let $\mathbf D$ be the diagonal matrix formed from the diagonal entries of $\mathbf A$. Jacobi uses $\mathbf M=\mathbf D$ and $\mathbf N=\mathbf D-\mathbf A$, giving:
  \[
    \mathbf G_J=\mathbf D^{-1}(\mathbf D-\mathbf A),
    \qquad
    \|\mathbf G_J\|_\infty
    =\max_i\frac{\sum_{j\ne i}|a_{ij}|}{|a_{ii}|}<1.
  \]
  The induced infinity-norm is the largest absolute row sum.
  For any eigenpair $(\lambda,\mathbf v)$ with $\mathbf v\ne\mathbf0$:
  \[
    \begin{aligned}
      |\lambda|\,\|\mathbf v\|_\infty
      &=\|\mathbf G_J\mathbf v\|_\infty
      \le\|\mathbf G_J\|_\infty\|\mathbf v\|_\infty,\\
      \rho(\mathbf G_J)&\le\|\mathbf G_J\|_\infty<1.
    \end{aligned}
  \]
  Thus, strict row diagonal dominance guarantees convergence of Jacobi.
  Proving convergence of Gauss–Seidel requires a separate argument.
  \lecturesource{Convergence: Diagonal dominance, Jacobi row-sum proof.}
\end{frame}

\begin{frame}{Rate of Convergence and an Error Bound}
  Suppose an induced matrix norm satisfies $q=\|\mathbf G\|<1$.
  Use its corresponding vector norm to bound the error:
  \[
    \begin{aligned}
      \|\mathbf e^{(k)}\|
      &=\|\mathbf G^k\mathbf e^{(0)}\|\\
      &\le\|\mathbf G^k\|\,\|\mathbf e^{(0)}\|\\
      &\le q^k\|\mathbf e^{(0)}\|.
    \end{aligned}
  \]
  The \hilite{contraction factor} $q$ satisfies
  $\|\mathbf e^{(k+1)}\|\le q\|\mathbf e^{(k)}\|$.
  A smaller $q$ gives a smaller upper bound for the same initial error.

  \vspace{0.7em}
  If the chosen norm gives $q\ge1$, the bound does not show that the error approaches zero,
  even when $\rho(\mathbf G)<1$.

  \vspace{0.5em}
  With $q<1$, each update reduces the error norm until the error is zero.
  \lecturesource{Convergence: Rate of convergence and an iteration bound.}
\end{frame}

\begin{frame}{A Sufficient Iteration Count}
  Let $E_0\ge\|\mathbf e^{(0)}\|$, with $E_0>0$. Choose a solution-error tolerance
  $0<\epsilon_e<E_0$ and assume $0<q<1$.
  Requiring $q^kE_0\le\epsilon_e$ gives:
  \[
    \begin{aligned}
      k\ln q&\le\ln(\epsilon_e/E_0),\\[3pt]
      k&\ge\left\lceil\frac{\ln(E_0/\epsilon_e)}{-\ln q}\right\rceil.
    \end{aligned}
  \]
  Dividing by $\ln q<0$ reverses the inequality. The ceiling rounds the count
  up to an integer. The error may reach the tolerance in fewer updates.

  \vspace{0.5em}
  For $q=0.5$, $E_0=1$, and $\epsilon_e=10^{-3}$, the bound gives
  \hilite{10 updates}, since $2^{-10}\approx9.77\times10^{-4}$.

  \vspace{0.5em}
  If $E_0\le\epsilon_e$, no update is needed. If $q=0$, one update gives the exact solution.
  \slidenote{Scope:}{$\epsilon_e$ bounds solution error; the algorithm uses residual tolerance $\epsilon$. The initial error bound $E_0$ may be unknown.}
  \lecturesource{Convergence: sufficient count, numerical anchor, and boundary cases.}
\end{frame}

\begin{frame}{Specific Methods}
  Split $\mathbf A$ into its diagonal matrix $\mathbf D$, strictly lower triangular
  part $\mathbf L$, and strictly upper triangular part $\mathbf U$:
  \[
    \mathbf A=\mathbf D+\mathbf L+\mathbf U.
  \]
  With nonzero diagonal entries and relaxation factor $\omega>0$,
  the matrices used in the correction solves are:
  \[
    \begin{aligned}
      \mathbf M_J&=\mathbf D, &&\text{Jacobi},\\[4pt]
      \mathbf M_{GS}&=\mathbf D+\mathbf L, &&\text{Gauss–Seidel},\\[4pt]
      \mathbf M_\omega&=\frac1\omega\mathbf D+\mathbf L,
        &&\text{successive over-relaxation}.
    \end{aligned}
  \]
  Each method converts the same residual into a different correction.
  \lecturesource{Specific Methods: common D + L + U convention and choices of M.}
\end{frame}

\begin{frame}{Jacobi: Independent Component Updates}
  For Jacobi, the correction solves $\mathbf D\mathbf d^{(k)}=\mathbf r^{(k)}$.
  Solving each row gives the corresponding component of the next iterate:
  \[
    x_i^{(k+1)}=\frac1{a_{ii}}
    \left(b_i-\sum_{j\ne i}a_{ij}x_j^{(k)}\right),
    \qquad i=1,\ldots,n.
  \]
  Every component update uses the \hilite{previous iterate} $\mathbf x^{(k)}$.
  Keep $\mathbf x^{(k)}$ unchanged until all components of $\mathbf x^{(k+1)}$ are ready.

  \vspace{0.6em}
  The iteration matrix is:
  \[
    \mathbf G_J=-\mathbf D^{-1}(\mathbf L+\mathbf U).
  \]
  Component updates can run in parallel. Strict row diagonal dominance
  guarantees convergence.
  \slidenote{Algorithm:}{\href{\jacobilink}{The Jacobi Method}}
  \lecturesource{Specific Methods: Jacobi. Component formula and iteration matrix from the Jacobi companion.}
\end{frame}

\begin{frame}{Gauss–Seidel: Sequential Component Updates}
  Solve $(\mathbf D+\mathbf L)\mathbf d^{(k)}=\mathbf r^{(k)}$
  by forward substitution. The equivalent component update is:
  \[
    x_i^{(k+1)}=\frac1{a_{ii}}\left(
      b_i-\sum_{j<i}a_{ij}x_j^{(k+1)}-\sum_{j>i}a_{ij}x_j^{(k)}
    \right).
  \]
  When updating component $i$, use the \hilite{new values} for components $j<i$
  and the values from the previous iterate for components $j>i$.
  The iteration matrix is:
  \[
    \mathbf G_{GS}=-(\mathbf D+\mathbf L)^{-1}\mathbf U.
  \]
  Strict row diagonal dominance guarantees convergence. Reusing updated values
  can reduce the iteration count, but Gauss–Seidel is not always faster than Jacobi.
  \slidenote{Algorithm:}{\href{\gslink}{The Gauss–Seidel Method}}
  \lecturesource{Specific Methods: Gauss–Seidel. Component formula and signed iteration matrix from the companion. Empty sums are zero.}
\end{frame}

\begin{frame}{Successive Over-Relaxation}
  The relaxation factor $\omega$ changes each sequential component update.
  The correction solves:
  \[
    (\mathbf D+\omega\mathbf L)\mathbf d^{(k)}=\omega\mathbf r^{(k)}.
  \]
  Applying the correction gives:
  \[
    x_i^{(k+1)}=(1-\omega)x_i^{(k)}
      +\frac{\omega}{a_{ii}}\left(
        b_i-\sum_{j<i}a_{ij}x_j^{(k+1)}-\sum_{j>i}a_{ij}x_j^{(k)}
      \right).
  \]
  When updating component $i$, the components $j<i$ already contain their relaxed values.
  The iteration matrix is:
  \[
    \mathbf G_\omega=(\mathbf D+\omega\mathbf L)^{-1}
      \bigl((1-\omega)\mathbf D-\omega\mathbf U\bigr).
  \]
  \slidenote{Algorithm:}{\href{\sorlink}{Successive Over-Relaxation}}
  \lecturesource{Specific Methods: SOR; correction, component formula, and iteration matrix from the SOR companion.}
\end{frame}

\begin{frame}{Relaxation and Convergence}
  The relaxation factor $\omega$ scales the change from the previous value
  of a component to the value obtained by solving its row equation:

  \vspace{0.5em}
  \begin{itemize}\setlength{\itemsep}{0.65em}
    \item $0<\omega<1$: \hilite{under-relaxation}, taking a partial step.
    \item $\omega=1$: the Gauss–Seidel update.
    \item $1<\omega<2$: \hilite{over-relaxation}, stepping beyond the row-equation solution.
  \end{itemize}

  \vspace{0.7em}
  If $\mathbf A$ is \hilite{symmetric positive definite}, convergence is
  guaranteed for $0<\omega<2$. Strict diagonal dominance alone does not
  guarantee convergence over that entire interval.

  \vspace{0.5em}
  A suitable factor can reduce the iteration count. The best choice depends
  on the matrix; formulas for an optimal factor need additional assumptions.
  \slidenote{Reading:}{\href{https://netlib.org/linalg/html_templates/node16.html}{Choosing the relaxation factor};
    \href{\sorlink}{special-case formula in the algorithm notebook}.}
  \lecturesource{Specific Methods: relaxation regimes and SPD guarantee.
    Netlib Templates, Choosing the Value of omega.
    SOR companion: symmetric positive definite tridiagonal special case.}
\end{frame}

\begin{frame}{Work per Iteration}
  A \hilite{sweep} updates all $n$ components once. For these methods,
  each completed iteration checks the residual and performs one full sweep.

  \vspace{0.6em}
  \renewcommand{\arraystretch}{1.45}
  \begin{tabularx}{\textwidth}{@{}lY@{}}
    \toprule
    Method & Correction work and dependencies\\
    \midrule
    Jacobi & Diagonal solve; components use the same previous iterate.\\
    Gauss–Seidel & Lower triangular solve; later rows depend on earlier updates.\\
    Relaxation & Relaxed lower triangular solve; the factor changes convergence.\\
    \bottomrule
  \end{tabularx}

  \vspace{0.7em}
  For dense matrices, a full iteration requires work proportional to $n^2$.
  Sparse implementations can work with the stored nonzero entries.

  \vspace{0.6em}
  Runtime depends on both the \hilite{cost per iteration} and the
  \hilite{number of iterations} needed to meet the residual tolerance.
  \lecturesource{Specific Methods: cost per sweep and total runtime.}
\end{frame}

\begin{frame}{Example: Fun with Iterative Methods}
  The example constructs a symmetric, strictly diagonally dominant matrix
  with positive diagonal entries. These properties make the matrix positive definite.

  \vspace{0.6em}
  \begin{enumerate}\setlength{\itemsep}{0.65em}
    \item Check symmetry and positive diagonal entries. Verify strict diagonal
      dominance in \hilite{every row}.
    \item Solve the system with an iterative method and a direct solver.
      Check the final residual norm and the norm of the difference between solutions.
    \item Benchmark both solvers with the same settings used above.
      Compare the timing distributions and memory allocations.
  \end{enumerate}

  \vspace{0.5em}
  The iterative solver forms explicit inverses and stores the iterates, adding to
  its runtime and memory use. The comparison applies to this implementation and test system.
  \slidenote{Example:}{\href{\examplelink}{Fun with Iterative Methods}}
  \lecturesource{Specific Methods: example callout and interpretation.
    Companion example: Tasks 1--3. Uses the corrected shared stopping loop;
    archived iterates stop at the first residual norm satisfying tolerance.}
\end{frame}

\begin{frame}{Lab: Compare Iterative Linear Solvers}
  In \href{\lablink}{L6d}, solve the same diagonally dominant system with
  Jacobi, Gauss–Seidel, and successive over-relaxation.

  \vspace{0.8em}
  \begin{itemize}\setlength{\itemsep}{0.8em}
    \item Compare each result with the \hilite{direct solution}.
    \item Use \hilite{residual-norm histories} and the solver's status flag to check
      whether the residual tolerance was met.
    \item Compare iteration counts and explore the effect of the
      relaxation factor.
  \end{itemize}

  \vspace{0.7em}
  Interpret speed only after checking the solution. A smaller iteration count
  is meaningful when the methods meet the same residual tolerance
  for the same system and initial guess.
  \slidenote{Lab:}{\href{\lablink}{Compare Jacobi, Gauss–Seidel, and SOR}}
  \lecturesource{Lab. Linked L6d notebook: learning objectives and Summary.}
\end{frame}

\begin{frame}{Summary}
  We developed stationary iterative methods by converting the residual
  into a correction through a matrix splitting.

  \vspace{0.4em}
  {\large\textbf{\ink{Key takeaways}}}
  \vspace{0.2em}
  \begin{itemize}\setlength{\itemsep}{0.65em}
    \item \textbf{Residual-based corrections.} The splitting determines the
      correction solve. Reaching an iteration limit does not establish convergence.
    \item \textbf{Convergence and error bounds.} The iteration matrix
      characterizes convergence. Diagonal dominance and contracting norms
      provide useful sufficient conditions and bounds.
    \item \textbf{Choosing a method.} Independent, sequential, and relaxed
      updates affect both computational work and convergence.
  \end{itemize}

  \vspace{0.6em}
  The example and lab connect these results to measured residual norms,
  iteration counts, and runtimes.
  \lecturesource{Summary: three key takeaways and concluding connection.}
\end{frame}

\end{document}
